\documentclass[a4paper]{article}

% Fandol lacks specialized CJK chars (e.g. 昇/騰); macOS → fontset=mac
\usepackage[fontset=mac]{ctex}
\usepackage{amsmath, amssymb}
\usepackage{graphicx}
\usepackage[margin=2.5cm]{geometry}
\usepackage[most]{tcolorbox}
\usepackage{etoolbox}
\usepackage{listings}
\usepackage{hyperref}
\usepackage{booktabs}
\usepackage{subcaption}
\usepackage{float}
\usepackage{tikz}
\usetikzlibrary{positioning}

% Highlight boxes
\newtcolorbox{knowledgebox}[1]{
    enhanced, colback=blue!5!white, colframe=blue!75!black, colbacktitle=blue!75!black,
    coltitle=white, fonttitle=\bfseries, title=#1, attach boxed title to top left={yshift=-2mm, xshift=2mm},
    boxrule=1pt, sharp corners
}

\newtcolorbox{importantbox}[1]{
    enhanced, colback=yellow!10!white, colframe=yellow!80!black, colbacktitle=yellow!80!black,
    coltitle=black, fonttitle=\bfseries, title=#1, sharp corners
}

\newtcolorbox{warningbox}[1]{
    enhanced, colback=red!5!white, colframe=red!75!black, colbacktitle=red!75!black,
    coltitle=white, fonttitle=\bfseries, title=#1, sharp corners
}

% Listings
\lstset{
    language=Python,
    basicstyle=\ttfamily\small,
    keywordstyle=\color{blue},
    stringstyle=\color{red!60!black},
    commentstyle=\color{green!60!black},
    breaklines=true,
    frame=single,
    numbers=left,
    numberstyle=\tiny\color{gray},
    captionpos=b,
    extendedchars=true,
    inputencoding=utf8
}

% Front-page metadata
\newcommand{\notetitle}{AI-Infra 规划课和模拟面——vLLM 大模型推理优化岗面试策略}
\newcommand{\noteauthors}{五道口纳什 \& Codex}
\newcommand{\notedate}{\today}
\newcommand{\videochannel}{五道口纳什}
\newcommand{\videopublishdate}{2026年4月18日}
\newcommand{\videoduration}{45分56秒}
\newcommand{\videourl}{https://www.bilibili.com/video/BV1UndWBHE8G/}
\newcommand{\videocoverpath}{cover.jpg}

\begin{document}

% ============================================================
% TITLE PAGE
% ============================================================
\begin{titlepage}
\centering
{\Large 课程笔记\par}
\vspace{1.2cm}
{\huge\bfseries \notetitle\par}
\vspace{0.8cm}
{\large \noteauthors\par}
\vspace{0.3cm}
{\large \notedate\par}
\vspace{1.2cm}

\ifx\videocoverpath\empty
{\small 请在生成笔记时填入视频封面图的本地路径。\par}
\else
\includegraphics[width=0.82\textwidth,height=0.45\textheight,keepaspectratio]{\videocoverpath}\par
\fi

\vfill
\begin{tcolorbox}[width=0.9\textwidth, colback=black!2!white, colframe=black!60, sharp corners]
\textbf{视频作者/频道}：\videochannel\par
\textbf{发布日期}：\videopublishdate\par
\textbf{视频时长}：\videoduration\par
\textbf{视频链接}：
\ifx\videourl\empty
未填写
\else
\href{\videourl}{\nolinkurl{\videourl}}
\fi
\end{tcolorbox}
\end{titlepage}

\tableofcontents
\newpage

% ============================================================
% SECTION 1: 行业格局
% ============================================================
\section{AI-Infra 行业格局与 vLLM 推理优化岗定位}

\subsection{AI-Infra 是什么}

AI-Infra（AI 基础设施）是连接模型算法与硬件的中间层，核心任务是在给定硬件上高效地训练和部署 AI 模型。随着大模型时代的到来，AI-Infra 已经从传统的「运维支持」角色演变为决定模型落地效率的关键环节。

AI-Infra 内部可细分为以下方向：

\begin{itemize}
  \item \textbf{训练框架}：分布式训练、混合精度、ZeRO 优化、Pipeline 并行。代表：Megatron-LM、DeepSpeed、PyTorch FSDP
  \item \textbf{推理框架/推理优化}：模型部署、KV Cache 管理、调度优化、量化部署。代表：vLLM、SGLang、TensorRT-LLM
  \item \textbf{算子开发}：CUDA Kernel 编写与优化、Attention 算子、GEMM 调优。代表：FlashAttention、cutlass
  \item \textbf{模型量化与轻量化}：GPTQ、AWQ、SmoothQuant、FP8/INT4 量化方案
  \item \textbf{芯片 SDK/Compiler}：将框架模型编译到专用硬件。代表：TensorRT、MLIR
\end{itemize}

\begin{figure}[H]
\centering
\begin{tikzpicture}[node distance=0.5cm, auto]
  \tikzstyle{layer}=[rectangle, draw, minimum width=12cm, minimum height=0.85cm,
                     align=center, font=\small, rounded corners=1pt]
  \node [layer, fill=blue!10] (app)    {应用层：ChatBot / CodeGen / 文生图 / Agent};
  \node [layer, fill=green!10, below of=app] (model)  {模型层：LLaMA / Qwen / DeepSeek / Stable Diffusion};
  \node [layer, fill=yellow!15, below of=model] (infra) {AI-Infra 层：训练框架 · 推理引擎 · 算子库 · 量化工具 · 芯片 SDK};
  \node [layer, fill=red!10, below of=infra] (hw)   {硬件层：NVIDIA GPU / 昇腾 NPU / AMD MI300X / 专用推理卡};
\end{tikzpicture}
\caption{AI-Infra 在技术栈中的定位}
\end{figure}

\subsection{推理优化岗的市场格局}

从岗位供需看，推理优化是当前 AI-Infra 方向中需求量最大的细分岗位。视频中提到上海一个城市就有约 100 个推理相关招聘岗位，且多数面向中小企业。

\begin{knowledgebox}{推理优化 vs 算子开发：岗位对比}
  \begin{itemize}
    \item \textbf{推理优化（推荐）}：大厂和互联网公司需求大，涉及全局优化思维（调度、KV Cache、PD 分离），发展路径宽
    \item \textbf{算子开发（谨慎选择）}：以国产芯片公司为主，需求集中但岗位总量偏少，薪资和天花板偏低。芯片公司对 PD 分离等业务场景优化无需求，它们提供的是 SDK 而非推理服务
  \end{itemize}
\end{knowledgebox}

\subsection{推理框架生态速览}

当前主流推理框架的技术定位差异：

\begin{enumerate}
  \item \textbf{vLLM}：社区最活跃的通用推理框架，PagedAttention + 连续批处理，生态完善
  \item \textbf{SGLang}：结构化生成见长，RadixAttention + 编程式前端，学术影响力大
  \item \textbf{TensorRT-LLM}：NVIDIA 官方方案，性能极致但灵活性受限
  \item \textbf{llama.cpp}：CPU/边缘推理标杆，量化方案丰富（GGUF）
\end{enumerate}

视频中讲者明确建议以 vLLM 作为学习入口，理由是社区活跃度高、岗位匹配面广、PD 分离等高级特性与业务场景强耦合。

\subsection{本章小结}

AI-Infra 的核心是「用系统手段解决模型落地效率」。推理优化岗目前需求最大；vLLM 是推荐的学习入口框架；算子开发需谨慎——它更多是芯片公司的需求，天花板和岗位总数都不如推理优化。

% ============================================================
% SECTION 2: Interview Strategy
% ============================================================
\section{vLLM 推理优化岗面试策略}

\subsection{讲者的面试方法论}

视频中讲者系统阐述了他面试候选人的完整策略，分为三个递进层次：

\begin{enumerate}
  \item \textbf{八股摸底（定位熟悉模块）}：先通过基础问答了解候选人对 vLLM 哪个模块最熟悉。这不是考核的全部，而是「定位入口」
  \item \textbf{模块深入（验证真实深度）}：针对最熟悉模块深入追问实现原理、设计决策、边界条件。如果候选人有社区 PR，则聚焦 PR 的工作展开——从发现问题到实现和优化的完整链路
  \item \textbf{发散性优化（考察系统思维）}：给一个业务场景（如「推理接口速度慢」）或开放式问题（如「如何优化长文本推理」），考察端到端的定位和优化能力
\end{enumerate}

\begin{importantbox}{面试核心逻辑}
  面试官不期望你每个模块都精通，但期望你在\textbf{一个模块}上有深入理解，并能展现\textbf{系统性的性能定位能力}。八股是入口而非终点；PR 是证明真实深度的最强信号。
\end{importantbox}

\begin{figure}[H]
\centering
\begin{tikzpicture}[node distance=2.2cm, auto]
  \tikzstyle{block}=[rectangle, draw, minimum width=4.5cm, minimum height=1.1cm,
                     align=center, font=\small, rounded corners=1pt]
  \tikzstyle{arrow}=[thick,->,>=stealth]
  \tikzstyle{branch}=[diamond, draw, minimum width=3cm, minimum height=1.1cm,
                      align=center, font=\small, aspect=2]
  
  \node [block, fill=blue!10] (start)  {八股摸底\\定位熟悉模块};
  \node [block, fill=green!10, right of=start, xshift=3cm] (deep)  {模块深入\\实现原理 / 设计决策};
  \node [block, fill=yellow!15, right of=deep, xshift=3cm] (open)  {发散性优化\\端到端定位 / 长文本优化};
  
  \draw [arrow] (start) -- node[above] {\small 确定深入方向} (deep);
  \draw [arrow] (deep) -- node[above] {\small 考察系统思维} (open);
\end{tikzpicture}
\caption{vLLM 推理优化岗面试三阶段}
\end{figure}

\subsection{PR 驱动面试的特殊路径}

对于社招转行或缺乏真实业务经验的同学，给目标框架提交社区 PR 是一条「破局路径」。讲者分享了真实案例：一位从大模型部署应用转行的候选人，因为给 vLLM 提交过两个 PR 而获得面试机会并最终通过。

PR 面试的问题链路：
\begin{itemize}
  \item 你是怎么发现这个问题的？（问题定位能力）
  \item 你的实现和优化思路是什么？（方案设计能力）
  \item 提交 PR 的过程是怎样的？（工程协作能力）
  \item 有没有被 reviewer 提过修改意见？你怎么处理的？（迭代能力）
\end{itemize}

讲者强调，PR 的含金量不在于代码行数多，而在于\textbf{能讲清楚「为什么做、怎么做、如何验证」}的完整 story。

\begin{warningbox}{PR 不是万能通行证}
  纯八股背得好但 PR 一问三不知 → 面试官会质疑 PR 真实性。PR 的价值在于它迫使你真正理解代码，而不是背面试题。
\end{warningbox}

\subsection{本章小结}

vLLM 面试的核心是三阶段递进：模块定位 → 深入验证 → 发散考察。PR 是证明能力的最强信号，但前提是你能把完整的问题-实现-优化链路讲清楚。发散性优化问题（如端到端性能定位、长文本推理优化）是区分候选人的关键。


% ============================================================
% SECTION 3: 项目经验与简历
% ============================================================
\section{项目经验与简历策略}

\subsection{玩具项目的致命伤}

视频中三位面试者均面临「玩具项目」困境。讲者指出，有经验的面试官能一眼识别学习性质的项目与真实业务项目：

\begin{enumerate}
  \item \textbf{缺少业务背景}：真实项目的优化一定有明确的业务诉求——「推理接口慢、用户投诉」或「某个模型上线后 P99 latency 不达标」。学习项目往往是「为了学 vLLM 而做 vLLM」
  \item \textbf{优化方向雷同}：P4S CUDA Graph、张量并行 + NCCL 通信——这些是课程标配，面试官每周能见到好几个写着同样关键词的简历
  \item \textbf{审美疲劳}：讲者提到 NanoVM 已经成为高频项目，「随便点开一个都是自己实现了一个推理框架」

\end{enumerate}

\begin{importantbox}{项目真实感 = 面试通过率}
  社招对项目经验的要求远高于校招。校招的容忍度稍高，但如果校招二面被问到「这是实验室项目还是自己做的」就尬住，说明项目包装的边界已经撞到了。
\end{importantbox}

\subsection{社区 PR 的含金量}

讲者多次强调，在缺乏真实业务经验的情况下，给 vLLM 社区提 PR 是最高含金量的项目加分项。原因：

\begin{itemize}
  \item \textbf{真实性无法造假}：PR 有完整的 review 记录、讨论历史、合入时间线
  \item \textbf{展示完整工程能力}：发现问题 → 设计方案 → 实现 → 测试 → review 迭代 → 合入
  \item \textbf{代码质量受社区检验}：能被 vLLM maintainer 接受的代码，本身就是能力背书
\end{itemize}

选择 PR 方向的建议：看 vLLM 的哪些特性可以与你已有的学习项目结合，在能力范围内选一个能实现的特性去做。

\subsection{项目包装的边界}

第三位面试者问「项目该怎么包装」，讲者的回答很明确：

\begin{itemize}
  \item 把 NanoVM 改名叫「自己实现的大模型推理框架」——没用。面试官关心的永远是\textbf{你做了什么、怎么做的、优化效果是什么}
  \item 把没做过的模块写进简历——更危险。写的越多，准备面越广，被问住的风险越大
  \item \textbf{关键是在现有项目中加入自己的优化}：比如 vLLM 有某个特性但 NanoVM 没实现，你去支持一下；或者你支持了一个 NanoVM 不支持的模型
\end{itemize}

\begin{warningbox}{简历膨胀的反噬效应}
  简历写「一整套大模型推理框架」意味着你需要准备调度器、KV Cache 管理、Continuous Batching、模型支持、量化方案……每个点都可能被深入问。宁可聚焦一个模块做深，也不要把自己摊薄。
\end{warningbox}

\subsection{芯片公司实习的「锁死效应」}

讲者提到一个值得警惕的职业路径现象：如果在芯片公司实习/工作过，后续找的还是一些芯片企业的实习，会被「一直卡在这里，很难跳去互联网厂商」。

原因：
\begin{itemize}
  \item 芯片公司的工作内容是 SDK 开发/算子适配，与互联网企业的推理优化岗（端到端性能定位、调度系统优化）的技能栈有差异
  \item 招聘方的惯性思维——「做过芯片公司的，可能不适应互联网的节奏和问题类型」
\end{itemize}

但这不意味着芯片公司绝对不能去——\textbf{关键是要问清楚工作内容}。如果是开发 TensorFlow 插件这种纯芯片适配工作，就不值得；如果是做推理框架本身的优化，那就和互联网岗位相通。

\subsection{本章小结}

项目经验的核心不是「项目名称是什么」，而是「你做了什么优化、怎么做的、效果如何」。社区 PR 是目前最强的项目证明。小心简历膨胀——写的越多，被问住的风险越大。芯片公司的选择需谨慎，关键是岗位工作内容而非公司招牌。

% ============================================================
% SECTION 4: 手撕与八股
% ============================================================
\section{手撕与八股准备}

\subsection{手撕的三类题型}

视频中明确区分了三种手撕题：

\begin{enumerate}
  \item \textbf{LeetCode 高频题}：算法和数据结构基础，面试标配，必须熟练
  \item \textbf{CUDA 算子手撕}：Reduce、MatMul（含 shared memory / bank conflict 优化版本）、LayerNorm 等常见算子。视频中第一位面试者的 MatMul 优化手撕被讲者建议作为重点准备
  \item \textbf{MHA（Multi-Head Attention）手撕}：从朴素的 PyTorch 版本到 CUDA Kernel 的完整实现路径
\end{enumerate}

CUDA 算子手撕的具体优化点（视频中提及的优化方向）：
\begin{itemize}
  \item Shared Memory 的使用（减少 global memory 访问）
  \item Bank Conflict 优化
  \item Warp-level 原语
  \item Tensor Core 利用（wmma / mma）
  \item FP16/BF16 精度处理
\end{itemize}

讲者推荐了一篇详细记录 MatMul 从基础 Naive 版本到高性能版本优化思路的文章，建议候选人「把它看熟练，面试时可以把里面的优化策略拿出来跟面试官交流」。

\begin{knowledgebox}{CUDA 算子手撕的准备策略}
  不是每个算子都需要精通手撕，但\textbf{一个经典算子（如 MatMul）的完整优化链路必须能讲清楚}。从 Naive → Shared Memory → Bank Conflict → Tensor Core，每一步的加速原理和瓶颈分析要烂熟于心。
\end{knowledgebox}

\subsection{八股的定位}

讲者澄清了一个常见误解：很多人以为框架面试就是背八股。实际上：

\begin{itemize}
  \item 八股只是第一个环节——用来\textbf{定位你熟悉 vLLM 的哪个模块}
  \item 定位之后，面试官会往深里问实现细节、设计权衡、边界情况
  \item 纯八股熟练但深入问题答不上来 → 面评会很差
\end{itemize}

所以八股是必要的起点，但不是终点。背诵 vLLM 的模块划分、调度流程、KV Cache 管理机制只是让你能坐到桌边——坐到桌边之后，要看你能不能真正聊出东西。

\subsection{本章小结}

手撕准备三块：LeetCode 高频题（保持手感）、CUDA 算子（至少精通一个算子的完整优化链路）、MHA 实现。八股是面试入口但非全部，核心是在熟悉模块上能深入阐述原理和设计决策。


% ============================================================
% SECTION 5: 学习路线与岗位选择
% ============================================================
\section{学习路线与岗位选择}

\subsection{推荐的入门路线}

讲者基于三个面试者的不同背景，给出了统一的学习路线建议：

\begin{enumerate}
  \item \textbf{CUDA 编程基础}：理解 GPU 架构（SM、Warp、Shared Memory、Global Memory）、掌握 CUDA C++ 基础、能手撕常见算子
  \item \textbf{vLLM 框架系统学习}：PagedAttention 原理、Scheduler 调度机制、Continuous Batching、KV Cache 管理、模型加载流程
  \item \textbf{选择一个深挖方向}：根据你的研究/项目背景，在 vLLM 的某个模块上达到面试深入水准
\end{enumerate}

\begin{figure}[H]
\centering
\begin{tikzpicture}[node distance=1.5cm, auto]
  \tikzstyle{phase}=[rectangle, draw, minimum width=6cm, minimum height=0.9cm,
                     align=center, font=\small, rounded corners=1pt]
  \tikzstyle{arrow}=[thick,->,>=stealth]
  
  \node [phase, fill=blue!10] (cuda)  {CUDA 编程基础\\GPU 架构 · 算子手撕 · 性能分析};
  \node [phase, fill=green!10, below of=cuda] (vllm)  {vLLM 框架系统学习\\PagedAttention · Scheduler · KV Cache};
  \node [phase, fill=yellow!15, below of=vllm] (deep)  {选择一个方向深挖\\调度优化 · PD 分离 · 量化部署 · 模型支持};
  \node [phase, fill=red!10, below of=deep] (pr)  {实战：社区 PR / 实习项目\\用真实代码和优化成果证明能力};
  
  \draw [arrow] (cuda) -- (vllm);
  \draw [arrow] (vllm) -- (deep);
  \draw [arrow] (deep) -- (pr);
\end{tikzpicture}
\caption{AI-Infra 推理优化方向学习路线}
\end{figure}

\subsection{CUDA 学到什么程度}

讲者对 CUDA 学习深度的定位很务实：
\begin{itemize}
  \item 目标是「面试时能把常见算子撕出来」，而不是「成为 CUDA 优化专家」
  \item 但如果面试的是算子开发岗，CUDA 的深度要求会更高
  \item 推荐方向：\textbf{推理优化岗}不需要 CUDA 达到算子专家水平，但必须理解 GPU 执行模型和基本优化方法
\end{itemize}

对于 PD 分离方向的同学，讲者特别指出 PD 分离必然和 vLLM/SGLang 联合使用，对框架的理解比 CUDA 深度更重要。

\subsection{Triton 要不要学}

第二位面试者问到 Triton，讲者的回答很明确：
\begin{itemize}
  \item \textbf{三四个月内找实习}：不用学 Triton。时间紧迫，先把 vLLM 框架熟悉起来
  \item \textbf{有充裕时间}：Triton 降低 CUDA 编程门槛，是值得掌握的技能
  \item 当前就业市场上直接要求 Triton 的岗位并不多
\end{itemize}

\subsection{算子开发 vs 推理优化：选哪个}

视频中最具决策价值的对比分析：

\begin{table}[H]
\centering
\begin{tabular}{p{3.5cm}p{5cm}p{5cm}}
\toprule
\textbf{维度} & \textbf{算子开发} & \textbf{推理优化} \\
\midrule
主要雇主 & 国产芯片公司 & 互联网大厂 + 中小厂 \\
岗位数量 & 较少 & 较多 \\
薪资水平 & 偏低 & 偏高 \\
工作内容 & CUDA Kernel 编写、适配芯片 & 端到端性能优化、调度系统 \\
业务关联 & 弱（不接触推理服务） & 强（直接服务业务） \\
职业跳槽 & 容易锁定芯片行业 & 互联网内流转灵活 \\
\bottomrule
\end{tabular}
\caption{算子开发 vs 推理优化岗位对比}
\end{table}

讲者的核心建议：\textbf{不是因为硬件限制才去做算子开发，而是因为你对算子本身感兴趣}。如果对算子开发「还好吧，也不是特别感兴趣」，就不要因为缺 GPU 而改方向——缺卡是暂时的，方向选择的影响是长期的。

\subsection{PD 分离与业务场景}

第二位面试者正在做 PD 分离相关研究（网络感知的调度优化），讲者给出了关键洞察：

\begin{itemize}
  \item PD 分离本身是一个有真实业务价值的优化方向——互联网企业部署推理服务时，Prefill 和 Decode 阶段的资源需求差异巨大
  \item 但「网络会议投递的 PD 分离论文」不等于「面试官眼中的 PD 分离经验」——关键在于你\textbf{是否理解 vLLM 中 PD 分离的实现和调度细节}
  \item 芯片公司不需要 PD 分离——因为芯片公司不部署推理服务，它们出售的是 SDK
\end{itemize}

\subsection{本章小结}

学习路线：CUDA → vLLM → 深挖一个方向 → 实战证明。CUDA 学到能过面试手撕即可，不必成为专家。Triton 短期内不必要。算子开发谨慎选择，推理优化岗位更优。PD 分离方向有价值但必须和 vLLM 框架深度结合。

% ============================================================
% SECTION 6: 三位面试者案例全景
% ============================================================
\section{三位面试者案例全景}

本视频的独特价值在于同时呈现了三个不同背景、不同阶段的面试者，形成了一幅完整的 AI-Infra 入行图景。

\subsection{案例一：社招转行者（6个月求职中）}

\textbf{背景}：已离职状态，之前有工作经验但非 AI-Infra 方向。每周 3-4 个面试机会，主要面中小企业。

\textbf{核心问题}：
\begin{itemize}
  \item 项目缺乏真实业务背景，面试官一眼能看出是学习项目
  \item 手撕准备不足——「感觉还不如群里研二实习生的手撕好」
  \item 虽然简历有 CUDA 优化相关内容（cutlass、算子优化），但面试时答得不够扎实
\end{itemize}

\textbf{讲者诊断}：
\begin{itemize}
  \item 简历本身写得不错，能拿到面试机会说明简历过关
  \item 核心差距在\textbf{项目真实感}和\textbf{手撕熟练度}
  \item 建议往 vLLM 社区提 PR——PR 是对项目真实性问题的最强回应
  \item 近期目标是进中小公司，积累真实经验后再考虑跳槽大厂
\end{itemize}

\textbf{关键启示}：社招的大厂门槛远高于校招，但不是没有路。先入行（中小公司）→ 积累真实项目经验 → 再跳大厂，是一条可行路径。

\begin{figure}[H]
\centering
\includegraphics[width=0.85\textwidth]{figures/scene_485.jpg}
\caption{第一位面试者（社招转行者）的模拟面场景 \protect\footnotemark}
\end{figure}
\footnotetext{视频画面时间区间：00:08:05--00:08:15。}

\subsection{案例二：研一学生（PD 分离研究方向）}

\textbf{背景}：研一，工作过一年后读研，导师指定 PD 分离方向，实验室资源有限（这学期无卡可用）。

\textbf{核心问题}：
\begin{itemize}
  \item 论文投网络会议还是系统会议？——怕网络会议不被 AI-Infra 面试官认可
  \item 无卡如何做框架研究？
  \item 要不要因为缺卡而转向算子开发 + 模型量化？
\end{itemize}

\textbf{讲者诊断}：
\begin{itemize}
  \item \textbf{继续投网络会议}——面试官关注的是你做了什么工作、解决了什么问题，不纠结会议类型。降级去投系统会议不划算
  \item \textbf{不要因为硬件限制改方向}——对算子「不是特别感兴趣」就不必勉强。后续实习时用公司的卡
  \item PD 分离的论文工作完全可以作为项目经验写进简历——「你做的工作、你的优化思路」才是面试官想看到的
  \item 必须学 vLLM——PD 分离在生产中必然和 vLLM/SGLang 联合使用，不懂框架会「挺尴尬的」
\end{itemize}

\textbf{关键启示}：研一就有方向意识是优势。硬件限制是暂时的（实习时有公司的卡），方向选择是长期的。PD 分离的方向和推理优化岗天然对口，不需要因为会议类型焦虑。

\begin{figure}[H]
\centering
\includegraphics[width=0.85\textwidth]{figures/scene_1440.jpg}
\caption{第二位面试者（研一 PD 分离方向）的模拟面场景 \protect\footnotemark}
\end{figure}
\footnotetext{视频画面时间区间：00:24:00--00:24:10。}

\subsection{案例三：研二学生（暑期实习求职中）}

\textbf{背景}：「二进宫」同学——今年 1 月咨询过一次，现在找暑期实习。一面能过，二面被淘汰。

\textbf{核心问题}：
\begin{itemize}
  \item 项目是 NanoVM + 简单优化（CUDA Graph 支持），面试官「审美疲劳」
  \item 二面被问到「这个项目是实习中的还是自己做的」就尬住
  \item 简历中的 Profiling 工作太常规，看不出差异化
  \item 在摩尔线程和无问苍穹两个实习 offer 之间犹豫
\end{itemize}

\textbf{讲者诊断}：
\begin{itemize}
  \item 关键不是包装——改个项目名字没有用。必须在项目里\textbf{加入自己的优化}：比如支持一个 NanoVM 没有的 vLLM 特性，或支持一个 NanoVM 不支持的模型
  \item 手上已有中小厂 offer 的同学：\textbf{先干着}。实习随时能跑路，积累真实经验比等不确定的池子更重要
  \item 摩尔线程 vs 无问苍穹：讲者倾向无问苍穹。理由是芯片公司的工作可能被锁定在芯片行业，且要问清楚工作内容——如果去摩尔线程是开发 TensorFlow 插件，那就不值得
  \item 无问苍穹虽然规模小（100+人），但「声量不错，算明星初创」，说明创始团队技术背景强
\end{itemize}

\textbf{关键启示}：简历上写「Profilng」不等于面试官能看到你的工作。你的优化必须\textbf{具体、可见、有量化效果}。实习先干着积累经验，不要空等不确定的机会。

\begin{figure}[H]
\centering
\includegraphics[width=0.85\textwidth]{figures/scene_2340.jpg}
\caption{第三位面试者（暑期实习求职）的模拟面场景 \protect\footnotemark}
\end{figure}
\footnotetext{视频画面时间区间：00:39:00--00:39:10。}

\subsection{本章小结}

三位面试者代表了 AI-Infra 入行的三种典型路径：社招转行（先入行再深耕）、研一规划（方向正确 + 时间充裕）、研二冲刺（差异化优化 + 先干再说）。共同教训是：项目真实感、手撕熟练度、方向选择，三者缺一不可。


% ============================================================
% SECTION 7: 总结与延伸
% ============================================================
\section{总结与延伸}

\subsection{讲者的结尾讨论}

视频末尾，讲者在三位面试者的对话之外给出了一些整体性的观察：

\begin{itemize}
  \item 中小公司对 AI-Infra 人才的需求正在快速增长——「上海大概有 100 个左右的推理招聘岗位」，且这个数字在持续上升
  \item 模型量化方向「吸收就业」的能力不如推理优化——量化方法趋于成熟（常用就那么几种），除非做量化算法本身的改进（偏算法岗）
  \item 面试中讲者最看重的是\textbf{系统思维}——不是能背多少八股，而是面对一个不熟悉的性能问题，能不能有条理地定位和优化
\end{itemize}

\subsection{核心主张与关键机制}

从本视频中蒸馏出的核心命题：

\begin{importantbox}{AI-Infra 入行三支柱}
  \begin{enumerate}
    \item \textbf{方向选择决定天花板}：推理优化 > 算子开发；互联网 > 芯片公司（除非你对芯片有特殊兴趣）
    \item \textbf{项目真实感是硬通货}：社区 PR > 实习项目 > 学习项目。真实感的本质是「能讲清楚为什么做、怎么做、效果如何」
    \item \textbf{系统思维是终极竞争力}：面试的高阶区分点不在八股背诵，而在于面对一个开放性问题（如「推理慢了怎么定位」）时的分析框架
  \end{enumerate}
\end{importantbox}

\subsection{实践含义}

\textbf{如果你在准备面试}：
\begin{itemize}
  \item 找一个 vLLM 的 issue，试着提一个 PR——哪怕只是修一个 bug 或加一个小 feature
  \item CUDA 算子手撕至少精通一个（MatMul 的完整优化链路是最安全的选择）
  \item 准备一个「端到端性能定位」的思维框架：profiling → 定位瓶颈 → 分析原因 → 设计优化方案 → 验证效果
\end{itemize}

\textbf{如果你是研一/研二学生}：
\begin{itemize}
  \item 尽早确定细分方向（推理优化 vs 算子开发 vs 量化），不要因为一时的硬件限制改方向
  \item 论文工作完全可以包装为项目经验——关键是展示你的分析和优化能力
  \item PD 分离、调度优化等方向与 vLLM 天然耦合，学框架是必须的
\end{itemize}

\textbf{如果你在考虑 offer}：
\begin{itemize}
  \item 问清楚岗位的工作内容——「推理框架优化」可能是做真正的推理引擎优化，也可能只是给芯片写 TensorFlow 插件
  \item 芯片公司经历可能锁定职业路径，但如果工作内容是推理框架本身（而非芯片适配），同样有价值
  \item 实习随时能跑路——先干着积累真实经验，不要空等不确定的机会
\end{itemize}

\subsection{开放问题}

视频没有覆盖但值得思考的方向：

\begin{itemize}
  \item \textbf{SGLang vs vLLM}：讲者提到了 SGLang 但也将其和 vLLM 并列。如果时间允许，是否应该两个框架都了解？SGLang 的 RadixAttention 和结构化生成在什么场景下优于 vLLM？
  \item \textbf{Attention 架构演进}：MLA（Multi-head Latent Attention，DeepSeek-V2/V3）对 KV Cache 管理的冲击——当 KV Cache 大小从 MB 级骤降到 KB 级，传统的 PagedAttention 是否还需要？
  \item \textbf{PD 分离的边界}：什么场景下 PD 分离是正收益？什么场景下反而是负优化（通信开销 > 并行收益）？这需要在实际业务中积累判断力
  \item \textbf{推理服务化趋势}：推理框架从库到服务的演进（vLLM 的 OpenAI-compatible API Server、SGLang 的 Runtime）对岗位技能要求的影响
\end{itemize}

\subsection{拓展阅读}

\begin{itemize}
  \item vLLM 官方文档：\href{https://docs.vllm.ai/}{https://docs.vllm.ai/} —— 从 PagedAttention 论文到最新特性的最佳入口
  \item \href{https://arxiv.org/abs/2309.06180}{Efficient Memory Management for Large Language Model Serving with PagedAttention} —— vLLM 核心论文
  \item \href{https://arxiv.org/abs/2312.07104}{SGLang: Efficient Execution of Structured Language Model Programs} —— SGLang 论文，了解 RadixAttention 和结构化生成
  \item CUDA C++ Programming Guide：\href{https://docs.nvidia.com/cuda/cuda-c-programming-guide/}{https://docs.nvidia.com/cuda/cuda-c-programming-guide/} —— CUDA 学习的官方基础
\end{itemize}

\subsection{本章小结}

AI-Infra 推理优化方向正处于需求爆发期。入行路径清晰：CUDA 打底 → vLLM 框架 → 深挖一个方向 → 实战证明。不同背景（社招转行/研一规划/研二冲刺）有各自的策略侧重，但共同的要求是：方向选择要慎重、项目要有真实感、系统思维是终极竞争力。先入行再深耕，是多数人最务实的路径。

\end{document}
